% !TEX root = paper.tex

\section{Introduction}

The Rust programming language~\rust has brought linear types into the mainstream.
Rust's sophisticated type system incorporates linear types and borrowing,
making it possible to write low-level systems code in a type-safe way
without requiring garbage collection.
This makes Rust an attractive language for developing low-level software
that needs both high performance and high assurance,
and Rust has gained rapid acceptance for system programming over the last
few years~\cite{rustandroid, rustlinux}.

Nevertheless, even type-safe code can still contain bugs
that harm a program's security and reliability.
Furthermore, systems programmers using Rust sometimes resort to unsafe code
(via Rust's \verb`unsafe` keyword)
for programming styles that do not fit into Rust's linearity discipline; e.g.,
it is awkward to encode doubly linked lists in Rust
because the backwards links violate linearity.

Formal verification promises to prove deeper properties about Rust programs,
including properties about low-level code that would otherwise require unsafe Rust features.
Hence, we introduce \verus, an SMT-based tool for verifying Rust code.
SMT-based verification can help
Rust overcome the limitations of Rust's strict typing discipline,
making it possible to safely express low-level code like doubly linked lists
or safe implementations of reader-writer locks for concurrent code.

Just as importantly,
we argue that Rust's linear type system can help make SMT-based verification easier,
bringing the power of substructural logics, like concurrent separation logic~\cite{DBLP:conf/lics/Reynolds02,DBLP:journals/tcs/OHearn07},
to SMT-based reasoning.
In particular, we demonstrate the use of linear ghost permissions
that enable a program to take specific actions on specific resources,
such as writing to a memory location.
Since the permissions are linear, they can track the evolving state of a resource
in the same way that separation logic formulas can track the state of a resource.
Since the permissions are ghost, they exist only during type checking and verification,
and do not impose any overhead on compiled executable code.

To take advantage of Rust's type system for checking linear ghost permissions,
\verus uses Rust to express specifications and proofs,
running Rust's linearity and borrow checking on the proofs.
By using a single language for specifications, proofs, and executable code,
\verus follows in the footsteps of earlier frameworks that combine proofs and programming,
such as Coq~\cite{coq}, Dafny~\cite{DBLP:conf/lpar/Leino10}, \fstar~\cite{mumon}, and Lean~\cite{lean2015}.
However, these systems were designed from scratch for verification,
unlike Rust, which was designed strictly as a programming language.
Therefore, it is important to ensure that the subset of Rust
that \verus allows in proofs and specifications is sound as a proof language.

In particular, we need to make sure that proofs terminate
and that functions used in specifications are pure, mathematical functions,
whereas executable code might contain infinite loops,
be nondeterministic, or have side effects.
In a language like Rust that contains recursive types, higher-order functions,
and type classes (traits), termination can be particularly subtle:
without positivity restrictions on recursive type definitions, for example,
these features together can encode nontermination.
Nevertheless, we want to restrict recursive types only for specifications and proofs,
not for executable code, unlike Coq, Dafny, and Lean,
which restrict all recursive types.

To enforce the distinction between specifications, proofs, and executable code,
\verus introduces a mode system that classifies all code as
specification, proof, or executable,
where specification and proof code are checked for termination.
All three modes of code are type checked;
proof and executable code are checked for linearity and borrowing;
only executable code is compiled to machine code.
We formalize this mode system using a small Rust-inspired lambda calculus,
proving preservation and progress for all well-typed expressions
and termination for all specification and proof expressions (\autoref{sec:formalization}).

\verus currently supports a large set of proof features
and a large subset of ordinary Rust features:
\begin{itemize}
\item Rust's finite range integers (i32, u32, etc.) as well as
  infinite-range integers (int, nat) for specifications and proofs
\item recursive algebraic datatypes (structs and enums),
  including mutual recursion, pattern matching, and pattern match guards
\item mutable variables, while loops, and return statements
\item recursive specification functions and inductive proofs,
  including mutual recursion and lexicographic decreases clauses
\item passing function arguments by borrowing (both \verb`&` and \verb`&mut`)
\item lifetime parameters on structs
\item generics (parametric polymorphism), with support for simple traits (type classes)
\item first-class functions (Rust closures) in specifications
\item modules with public and private definitions
\item preconditions, postconditions, and loop invariants
\item quantifiers (forall, exists, choose) with both automated and manual SMT trigger selection,
  as well as integrated quantifier profiling to diagnose verification performance issues
\item programmer control of SMT performance by selectively hiding and revealing
  specification function definitions,
  including control over recursive function unrolling
\item support for bit-vector reasoning and proof by computation
\item strings and characters
\item a library of types for specifications, including sequences, sets, and maps
\item low-level pointer reasoning
\item concurrency and state machines
\end{itemize}

Verus is also able to handle some situations that require \texttt{unsafe}.
For example, Verus is able to verify the use of raw pointers and unsafe cells, which can be useful
for some low-level pointer reasoning, lock implementations, and interior mutability use-cases.
As we will see, \verus' support for linear ghost state is crucial for this support.
Note though, that Verus does not attempt a full aliasing and provenance model for Rust's pointers;
our simplified model for ``raw pointers'' only handles those that point into the global heap.



Some features are still missing, notably support for separate verification of multiple crates
and functions that return mutable references.
However, we believe that supporting crates and various other Rust features is a matter of engineering,
and supporting functions that return mutable references can follow earlier research (\autoref{sec:related})
by Prusti~\cite{DBLP:conf/nfm/AstrauskasBFGMM22},
RustHorn~\cite{DBLP:conf/esop/0002T020},
Creusot~\cite{denis:hal-03737878},
and Aeneas~\cite{DBLP:journals/pacmpl/HoP22}.


Regardless, this paper will not focus on all the features supported by \verus,
but will instead focus on the most novel contributions:
\begin{enumerate}
\item usage of Rust's linearity and borrow checking in proofs
\item verification of pointer-manipulating Rust code and concurrent Rust code,
  based on a combination of linearity, borrowing, and SMT solving
\item a mode system for enforcing the different properties of specs, proofs,
  and executable code
\item formalization of the mode system, including checking of linearity and borrowing
\end{enumerate}

The remainder of this paper
introduces \verus by example (\autoref{sec:design});
discusses handling unsafe code (\autoref{sec:unsafe});
applies \verus to pointer-based code (including a doubly linked list, \autoref{sec:linear-ghost}),
interior mutability (\autoref{sec:interior-mutability}),
and concurrent code (\autoref{sec:concurrency});
discusses Verus' implementation (\autoref{sec:implementation}),
user experience (\autoref{sec:user-experience}), and limitations (\autoref{sec:limitations});
and presents syntax, semantics, and proofs for a formal lambda calculus
with modes, linearity, and borrowing (\autoref{sec:formalization}).
